\section{Síntesis y ampliación}

El Model-Based RL es un componente importante de la caja de herramientas del RL; su principal ventaja es la eficiencia en el uso de datos.

Puntos principales:
\begin{enumerate}
    \item Aprender un modelo de la dinámica puede reducir mucho los datos de interacción real que se necesitan
    \item El modelo puede usarse para planning (CEM/MPC) o para generar datos sintéticos (al estilo de Dyna)
    \item El error del modelo es el desafío central; se mitiga con ensambles, rollouts de horizonte corto y actualización continua
    \item Es especialmente valioso en escenarios donde los datos son costosos, como la manipulación robótica
\end{enumerate}

\begin{knowledgebox}{Lecturas adicionales}
\begin{itemize}
    \item Nagabandi et al., ``Neural Network Dynamics for Model-Based Deep RL with Model-Free Fine-Tuning,'' ICRA 2018
    \item Chua et al., ``Deep Reinforcement Learning in a Handful of Trials using Probabilistic Dynamics Models (PETS),'' NeurIPS 2018
    \item Janner et al., ``When to Trust Your Model: Model-Based Policy Optimization (MBPO),'' NeurIPS 2019
    \item Sutton, ``Dyna, an Integrated Architecture for Learning, Planning, and Reacting,'' 1991
\end{itemize}
\end{knowledgebox}

\end{document}
